% !TEX program = xelatex
% 공통수학2 · Ⅰ. 도형의 방정식 · 01 선분의 내분 · 3/3차시
% 교사용 지도안

\documentclass[11pt,a4paper]{article}

\usepackage{kotex}
\usepackage{iftex}
\ifPDFTeX\else
  \setmainhangulfont{Noto Serif CJK KR}
  \setsanshangulfont{Noto Sans CJK KR}
\fi

\usepackage[margin=2.1cm]{geometry}
\usepackage{parskip}
\usepackage{enumitem}
\usepackage{array}
\usepackage{tabularx}
\usepackage{booktabs}
\usepackage{xcolor}
\usepackage{amssymb}
\usepackage{tikz}
\usepackage[most]{tcolorbox}
\usepackage{fancyhdr}
\usepackage{titlesec}

\definecolor{goldc}{HTML}{B07C22}
\definecolor{goldstrong}{HTML}{8A5F16}
\definecolor{indigoc}{HTML}{2B3B57}
\definecolor{indigosoft}{HTML}{6E82A6}
\definecolor{linec}{HTML}{D6DACB}
\definecolor{surface2}{HTML}{E4E8DC}
\definecolor{notebg}{HTML}{FBF3E3}
\definecolor{goodbg}{HTML}{E7EEE0}
\definecolor{inksoft}{HTML}{52604F}

\titleformat{\section}
  {\normalfont\sffamily\bfseries\large\color{indigoc}}
  {\thesection}{0.6em}{}
\titlespacing{\section}{0pt}{16pt}{8pt}

\newtcolorbox{ideabox}{
  enhanced, breakable, colback=white, colframe=linec, boxrule=0.5pt,
  leftrule=3pt, colframe=goldc, arc=2pt,
  left=10pt, right=10pt, top=8pt, bottom=8pt
}
\newtcolorbox{calloutbox}{
  enhanced, breakable, colback=white, colframe=indigosoft, boxrule=0.5pt,
  leftrule=3pt, arc=2pt, left=10pt, right=10pt, top=7pt, bottom=7pt
}
\newtcolorbox{notebox}{
  enhanced, breakable, colback=notebg, colframe=goldc, boxrule=0.5pt,
  leftrule=3pt, arc=2pt, left=10pt, right=10pt, top=7pt, bottom=7pt
}
\newtcolorbox{answerbox}{
  enhanced, breakable, colback=white, colframe=linec, boxrule=0.5pt,
  arc=2pt, left=10pt, right=10pt, top=7pt, bottom=7pt
}

\newcommand{\lessonbanner}[3]{%
  \begin{tcolorbox}[enhanced, colback=indigoc, colframe=indigoc,
    boxrule=0pt, arc=0pt, left=14pt, right=14pt, top=14pt, bottom=10pt,
    borderline south={2.4pt}{0pt}{goldc},
    fontupper=\color{white}]
    {\sffamily\footnotesize\color{white!85!black} #1}\\[4pt]
    {\sffamily\bfseries\Large #2}\\[3pt]
    {\sffamily\small\color{white!88!black} #3}
  \end{tcolorbox}
}

\pagestyle{fancy}
\fancyhf{}
\renewcommand{\headrulewidth}{0pt}
\fancyfoot[L]{\footnotesize\color{inksoft} 공통수학2 · 2026학년도 2학기 · 지평선고등학교}
\fancyfoot[R]{\footnotesize\color{inksoft} 교사용 지도안 -- 배포 금지}

\begin{document}

\lessonbanner{공통수학2 · Ⅰ. 도형의 방정식 · 1. 평면좌표와 직선의 방정식}%
  {01 선분의 내분 --- 3/3차시}%
  {무게중심과 황금비: 탐구\,\&\,융합 · 교사용 지도안 (2026학년도 2학기)}

\vspace{10pt}

\noindent
\begin{tabularx}{\linewidth}{@{}>{\bfseries\color{indigoc}}p{2.6cm}X@{}}
\toprule
대상 & 고1 · 2개 반 · 반당 13명 \\
차시 & 50분 (도입 10 · 전개 30 · 정리 10) \\
성취기준 & 내분점의 좌표를 구하는 공식을 활용하여 삼각형의 무게중심의 좌표를 구하고, 여러 가지 문제를 해결할 수 있다. \\
선행 학습 & 1$\sim$2차시 --- 수직선·좌표평면 위의 선분의 내분 \\
\bottomrule
\end{tabularx}

\section{학습 목표 \& Big Idea}

\begin{ideabox}
\textbf{\color{goldstrong}영속적 이해 (Big Idea)}\\
`비례로 위치를 계산한다'는 원리는 두 점을 넘어 세 점, 나아가 도형 전체의 `균형점'을 찾는 데까지 확장된다. 같은 도구가 다른 크기의 문제를 푼다.
\vspace{4pt}

\small\color{inksoft}
핵심개념: \textbf{균형(평균)} \quad\textbullet\quad 핵심개념: \textbf{비례의 재귀적 적용} \quad\textbullet\quad
일반화: \textbf{같은 원리가 규모를 바꿔 가며 반복 적용된다}
\end{ideabox}

\begin{calloutbox}
\textbf{차시 학습 목표}
\begin{itemize}[leftmargin=1.4em, itemsep=2pt, topsep=4pt]
  \item 삼각형의 세 중선이 만나는 점(무게중심)의 좌표를 내분점 공식으로 유도하고 구할 수 있다.
  \item 도형을 변형해도 무게중심이라는 `균형의 성질'이 보존되는 경우를 탐구할 수 있다.
  \item 내분의 개념이 황금비라는 미적 비례와 연결됨을 이해하고, 주변 사례를 탐구할 수 있다.
\end{itemize}
\end{calloutbox}

\section{질문의 연결}

\begin{tabularx}{\linewidth}{@{}>{\bfseries\color{indigoc}\small}p{2.4cm}X@{}}
사실적 질문 & 삼각형의 무게중심의 좌표는 어떻게 구할까? \\[6pt]
개념적 질문 & 왜 세 꼭짓점 좌표의 평균이 곧 `균형점'이 될까? \\[6pt]
논쟁적 질문 & 수학적으로 `아름다운 비율'(황금비)이 존재한다는 생각에 동의하는가? 아름다움을 숫자로 정의할 수 있을까? \\
\end{tabularx}

\section{수업 흐름}

\begin{tabularx}{\linewidth}{@{}>{\centering\arraybackslash}p{1.6cm}X>{\raggedright\arraybackslash\small}p{3.1cm}@{}}
\toprule
\textbf{단계} & \textbf{활동 \& 발문} & \textbf{ATL / VTR} \\
\midrule
\textcolor{goldstrong}{\textbf{도입}}\newline\footnotesize 10분 &
\textbf{마음공부 · 유무념 대조 (2분)} --- ``이번 주 나의 하루 중 `균형'이 잘 맞았던 순간은?''\newline
\textbf{생각 열기 (8분)} --- 삼각형 모양 판자를 손가락 하나로 받쳐 균형을 잡는 사진/영상 제시 $\to$ ``어느 점에서 받쳐야 안 쏟아질까?'' 직관적 예측 후 논의. &
ATL: 사고(예측)\newline VTR: See-Think-Wonder \\
\addlinespace
\textcolor{goldstrong}{\textbf{전개}}\newline\footnotesize 30분 &
\textbf{개념 형성 (10분)} --- 중선의 정의(꼭짓점과 대변 중점을 잇는 선분) 확인 $\to$ 변 BC의 중점 M 계산 $\to$ 중선 AM을 $2:1$로 내분하는 점이 무게중심임을 안내 $\to$ 내분점 공식 적용으로 $G=\left(\dfrac{x_1+x_2+x_3}{3},\dfrac{y_1+y_2+y_3}{3}\right)$ 직접 유도(교사가 답을 주지 않고 학생이 계산).\newline
\textbf{탐구 활동 (12분)} --- 삼각형 A$(6,3)$, B$(0,-3)$, C$(-3,9)$의 무게중심을 구한 뒤, 세 변 AB, BC, CA를 각각 $2:1$로 내분한 점 D, E, F로 만든 새 삼각형의 무게중심을 다시 구해 두 값이 일치하는지 확인 (모둠 활동, 결과를 칠판에 취합). --- 지도서 원문 문제 6.\newline
\textbf{탐구\&융합 (8분)} --- 황금비 소개: $\overline{AP}:\overline{PB}=\overline{AB}:\overline{AP}$를 만족하는 황금분할점. 명함·신용카드의 가로세로 비를 직접 측정해 황금비($\approx 1.618$)에 가까운지 확인. &
ATT: 촉진자 역할\newline ATL: 협업·조사\newline VTR: Connect-Extend-Challenge \\
\addlinespace
\textcolor{goldstrong}{\textbf{정리}}\newline\footnotesize 10분 &
\textbf{모둠 결과 공유 (5분)} --- 두 무게중심이 일치한다는 결과를 모둠별로 발표 $\to$ ``왜 일치할 수밖에 없는가''를 논리적으로 설명해 보기.\newline
\textbf{성찰 (5분)} --- 감사 일기 1문장 + 소단원(선분의 내분) 마무리 소감 + 다음 차시 예고(02 직선의 위치 관계). &
마음공부 · 감사 일기 \\
\bottomrule
\end{tabularx}

\vspace{8pt}
\begin{minipage}[t]{0.48\linewidth}
\begin{calloutbox}
\textbf{지도상의 유의점}
\begin{itemize}[leftmargin=1.2em, itemsep=2pt, topsep=4pt]
  \item 무게중심 공식을 곧바로 암기시키지 말고, 반드시 중점 $\to$ 중선의 $2:1$ 내분이라는 2단계 유도를 학생이 직접 계산하게 한다.
  \item 탐구 활동의 두 삼각형 무게중심이 같다는 결과는 우연이 아니라 벡터적으로 증명 가능함을 짚어 주되(고등 수준을 넘는 증명은 생략), ``항상 성립한다''는 확신을 주는 것이 핵심이다.
  \item 황금비 활동은 정밀 측정보다 `비슷한 비율을 스스로 찾아보는 경험'에 방점을 둔다.
\end{itemize}
\end{calloutbox}
\end{minipage}\hfill
\begin{minipage}[t]{0.48\linewidth}
\begin{notebox}
\textbf{인문학적 탐구 연결}\\
무게중심은 ``전체를 대표하는 하나의 점''이라는 점에서, 한 공동체나 팀의 의견이 어떻게 `균형점'을 찾아가는지에 대한 은유로 확장할 수 있다. 황금비는 수학과 예술·건축이 만나는 지점으로, 다음 소단원(직선의 위치 관계)으로 넘어가기 전 이 소단원 전체를 인문학적으로 갈무리하는 역할을 한다.
\end{notebox}
\end{minipage}

\section{탐구 활동 답안}

\begin{answerbox}
삼각형 A$(6,3)$, B$(0,-3)$, C$(-3,9)$의 무게중심: $G=\left(\dfrac{6+0-3}{3},\dfrac{3-3+9}{3}\right)=(1,3)$

\vspace{4pt}
변 AB, BC, CA를 각각 $2:1$로 내분한 점: D$(2,-1)$, E$(-2,5)$, F$(3,5)$

\vspace{4pt}
삼각형 DEF의 무게중심: $\left(\dfrac{2-2+3}{3},\dfrac{-1+5+5}{3}\right)=(1,3)$

\textbf{\color{goldstrong}$\to$ 두 무게중심이 $(1,3)$으로 정확히 일치한다.} (지도서 원문 문제 6)
\end{answerbox}

\section{성취수준 (이 차시 범위)}

\begin{tabularx}{\linewidth}{@{}>{\bfseries\color{goldstrong}}p{0.9cm}X@{}}
\toprule
A & 무게중심 공식을 유도 과정과 함께 설명하고, 변형된 도형에서도 무게중심이 보존되는 이유를 논리적으로 설명할 수 있다. \\
B & 무게중심 공식을 이해하고 정확히 계산하며, 탐구 활동의 결과를 설명할 수 있다. \\
C & 공식을 이용해 삼각형의 무게중심을 계산할 수 있다. \\
D & 안내에 따라 무게중심을 계산할 수 있다. \\
E & 세 점의 $x$좌표, $y$좌표의 평균을 각각 계산할 수 있다. \\
\bottomrule
\end{tabularx}

\section{다음 차시}

\begin{calloutbox}
\textbf{4/4차시 --- 02 직선의 위치 관계 (1/2).} `01 선분의 내분' 소단원이 여기서 마무리된다. 다음 시간부터는 두 직선의 평행 조건과 수직 조건을 탐구한다. 오늘 다룬 `비례'의 언어가 이번엔 `기울기'라는 새로운 도구로 이어짐을 도입에서 짚어 줄 것.
\end{calloutbox}

\end{document}
